\documentclass[10pt]{book}

\usepackage{fourier}
\usepackage{fancyhdr}
\usepackage{color}
\usepackage{graphicx}
\usepackage{soul}

\pdfpageheight 8in
\pdfpagewidth 5in

\setlength\topmargin{-.344in}
\setlength\headheight{0in}
\setlength\headsep{0in}
\setlength\textheight{6in}
\setlength\textwidth{3.2472in}
\setlength\oddsidemargin{-.25in} % to make effective .5in b/c of .25 of binding
\setlength\evensidemargin{0in} % to make effective .5in b/c of .25in of binding
\setlength\parindent{0in}
\setlength\parskip{1em}

\newenvironment{fullpage}{%
  \begin{list}{}{%
    \setlength{\topmargin}{-1in}
    \setlength{\textheight}{8in}
    \setlength{\topsep}{0pt}%
    \setlength{\leftmargin}{-.6in}% take into account the .25" of binding
    \setlength{\listparindent}{\parindent}%
    \setlength{\itemindent}{\parindent}%
    \setlength{\parsep}{\parskip}%
  }%
  \item[]}{\end{list}}

\def\ck{\kern 1pt}
\def\phd{{\sc PhD}}

\sodef\chpso{\Large}{.1em}{.5em}{0em}
\sodef\secso{}{.05em}{.3em}{0em}

\makeatletter
\def\cleardoublepage{\clearpage\ifodd\c@page\else\thispagestyle{empty}\hbox{}\newpage\fi}

\def\@chapter[#1]#2{\cleardoublepage%
\refstepcounter{chapter}\addcontentsline{toc}{chapter}{#1}%
\chpso{\MakeUppercase{#1}}\\*
\rule[-2pt]{\linewidth}{.5pt}\vspace*{-9pt}}

\renewcommand{\section}{\@startsection{section}{1}{0mm}%
   {10pt}%
   {8pt}{\large\scshape\MakeLowercase}}

\renewcommand{\subsection}{\@startsection{subsection}{2}{0mm}%
   {6pt}%
   {6pt}{\normalfont\normalsize\itshape}}

\renewcommand{\subsubsection}{\@startsection{subsubsection}{3}{0mm}%
   {0pt}%
   {0pt}{\normalfont\normalsize\scshape\MakeLowercase}}

\renewcommand{\tableofcontents}{\cleardoublepage
\chpso{TABLE OF CONTENTS}
  \vskip0pt
  \rule[8pt]{\linewidth}{.5pt}
  \vspace{-3pt}
  \@starttoc{toc}}
\makeatother

\pagestyle{fancy}
\fancyhf{}
\fancyfoot[LE,RO]{\addvspace{10pt}\small\thepage}
\renewcommand{\headrulewidth}{0pt}
\fancypagestyle{plain}{%
\fancyhead{}
}

\begin{document}

\frontmatter

% 1st right page: blank
\thispagestyle{empty}
\mbox{}
\newpage

% 1st left page: blank
\thispagestyle{empty}
\mbox{}
\newpage

% 1st right page: promo
\thispagestyle{empty}
\mbox{}
\vfil\eject

% 2nd left page: blank
\thispagestyle{empty}
\mbox{}
\newpage

% 2nd right page: half-title
\thispagestyle{empty}
\begin{fullpage}
\includegraphics{2007-redbook-halftitlepage.pdf}
\end{fullpage}
\vfil\eject

% 3rd left page: blank
\thispagestyle{empty}
\mbox{}
\newpage

% 3rd right page: full-title
\thispagestyle{empty}
\begin{fullpage}
\includegraphics{2007-redbook-titlepage.pdf}
\end{fullpage}
\vfil\eject

% 4th left page: copyright, publication information
\vspace*{5in}
{\small
\noindent Published 2007\\*
Printed by CafePress.com in the United States of America

\noindent Typeset with \LaTeX\ in Utopia, originally by Austen McDonald\\*
http://austenmcdonald.com
}
\thispagestyle{empty}
\vfil\eject

% 4th right page: dedication
\thispagestyle{empty}
{\Large\it Dedicated to...}
\vskip 0pt
{\it\hspace{.5in} the incoming class of 2007...best of luck!}
\newpage

% 5th left page: blank
\mbox{}
\thispagestyle{empty}
\newpage

% 5th right page: acknolwedgements: put all authors and contributors here, in
% alphabetical order, grouped by year
\chpso{ACKNOWLEDGMENTS}
\vskip0pt
\rule[8pt]{\linewidth}{.5pt}
\vspace{-3pt}
{\bf 2007}\\*
Ioannis Antonellis\\*
Billy Chen\\*
Mike Houston\\*
Jeff Klingner\\*
Austen McDonald\\*
Adam Oliner\\*
Doantam Phan\\*
Ron Yeh\\*

{
\setlength\parskip{0em}
\tableofcontents
}

\mainmatter

\section{\secso{What is This?}}
The Redbook is the attempt of a few \phd\ students to pass on bits of wisdom
and frank advice to newer students in the hope that the next generation avoids
their missteps. It lives online at http://theredbook.pbwiki.com.

\chapter{The Basics}

\section{\secso{Three Things to Know}}
\begin{enumerate}
\item Academia is a business, and graduate student is a job title. From Azuma:
``Academia is not the Real World and does not work in the same way that the
ordinary corporate world does. However, it is a business nonetheless and as a
graduate student, you must treat it that way. Graduate school made a lot more
sense and became much easier for me after I realized this. If you think of
graduate school as an `Ivory Tower' free of politics, money problems, and
real-world concerns, you will be severely disappointed.''
\item Remember you're here to get a \phd. Especially if you come straight from
college, the apparent lack of requirements may cause you to not to take school
seriously. We're not saying don't have fun, just don't have {\it too much fun}. Have
a plan about how you're going to complete your degree in a timely fashion.

\item Good advice can be gotten by chatting with senior graduate students behind
closed doors, as they will be much more willing to be open and honest.

\end{enumerate}

\section{\secso{Your First Quarter}}
What should I be doing in my first few months at Stanford?

In your first quarter, {\it don't drift}. You should be studying for comps and
finding an advisor. You may hear that comps are ``trivial'' and you should not
waste time studying for them because you're so smart you'll just pass them
without effort. This may be true for some, but it's our opinion that you should
study hard for comps. Forming study groups with other new students is a
great idea. You may want to take a class the first quarter. Finding an
advisor won't take all your time, and taking a class will help you meet new people.
If you have weak comp areas, taking a class or two can help you pass out of them.

Getting comps over earlier is better. As time passes, you'll have commitments
to research projects and paper deadlines. At that point, studying for comps or
taking classes to pass them can take time away from what's important.

If you're not sure what area you want to be in, work with a professor for a
trial quarter. Don't just attend meetings: working on an existing project with
more senior grad students is a great opportunity to get mentored and learn by
their example. (Try to pick a good project!). Take a lesson from other fields,
which have students rotate through different labs in their first year. Even if
you don't end up with that advisor, you've gotten research experience,
perhaps credit on a paper, and a professor in the department who knows you. If you have
difficulties later in the program, for whatever reason, having a
professor who knows you and who will go to bat for you is key.

\section{\secso{Advisors}}
\subsection{Choosing an Advisor}
{\it The choice of advisor is by far the most important decision you'll make} (by an
order of magnitude). 

It will completely control your grad school experience and
will affect your post-grad-school career as well, especially if you want to
continue in academia.

\subsubsection{Other Students' Experiences. }
The best way to find out what an advisor is like to is to chat with current
members of that professor's group behind closed doors. They'll give you the
real information about a professor, from inspiring tales to horror stories.
Both exist. Ask them about the topics below:

\subsubsection{Funding. }
When choosing an advisor, you need to understand their funding situation. Don't
be afraid to discuss this candidly---it's the only way to avoid
misunderstandings. Without funding, you will be supported by the department,
which means you have to be a course assistant ({\sc ca}). You don't want this
to happen because {\sc ca}ships take a considerable amount of time away from
research. A 50\% {\sc ca}ship covers your tuition and stipend, but generally
requires $\approx 20$ hours per week of work from you. Although 25\% {\sc ca}ships
exist, the common wisdom is that you will be working almost as much as the 50\%
{\sc ca} for less money. Other sources of funding include fellowships
(Microsoft, {\sc ibm}, {\sc nvidia}, {\sc ati}) but usually the department controls
who gets to apply for these fellowships (or it is limited to one applicant per
professor).

Questions to keep in mind: Are there active grants that can
support your specific research project as a 50\% {\sc ra}? How many quarters of
support are these grants? Are they renewable? What is the criteria under which
they can be renewed? Are your professors actively writing more grants for the
future? If not, will they help you write a grant (it's good practice, in any
case)?

It's bad to be forced into {\sc ca}ing quarter after quarter,
trying to produce enough research to keep your advisor happy. Have a candid
talk with your advisor so this does not happen to you. However, if this
does happen to you, and you still end up graduating, a good way to turn those
lemons into lemonade is to make sure you graduate with a Distinction in
Teaching.

\subsubsection{Advisor Age. }
A young (i.e., pre-tenure) advisor will be hungry for results and will be
closely involved with your work, which is a good thing. They are likely to
closely monitor the relevant conferences and will be
well aware of related work. The downside is that they may fail to make tenure,
which could screw up your grad school career.

An old advisor who has already made their mark has lots of
experience and contacts. They are a safer bet than a pre-tenure professor since
they aren't going anywhere, and many have reliable funding. However, a
common experience among students of older, already successful professors is
that they aren't hungry for success anymore, and will be only superficially
engaged in the research of their students, leaving the students feeling adrift.

\subsubsection{Recent Publishing Record. }
If you want to pursue an academic career at a top-tier school, you'll need
several high-quality first-author papers. Look at how many papers the
professor's current students have published in the past few years, and at what
conferences. This is an indicator of where the professor
stands with respect to the earlier point: are they still hungry to achieve
success, or are they resting on their laurels?

Aside from this question of hunger, professors are also just different from
each other with respect to the way they approach conferences; some publish
papers three times more often than others. Of course, paper quality matters
more than quantity, so use Citeseer, ACM Digital Library, or Google Scholar to see how often the
professor's recent papers are cited, as an indicator of how well their work is
received by their research community.

\subsubsection{Advising Style. }
Many will suggest you ask if a professor is ``hands-on'' or ``hands-off.'' Unless
you lean strongly toward one or the other, stick to other metrics when picking
an advisor. Be wary of professors at either extreme (talk to senior grad students).

\subsubsection{Group Composition. }
You'll want an advisor with several current students, since their body of
knowledge and experience will be of great benefit. Many groups work a little
like an apprenticeship: young students join existing projects to learn from the
older students, and as they get more experience they will strike out in their
own directions.

If all the current students are near graduation and there are no younger
current students in the pipeline, this could be a problem.

Be wary of those advisors with really large groups (15+ students), since you'll
get much less face-time than with an advisor with only 5--10 students.

\subsubsection{Industry Contacts. }

If you are in a systems-related field, then industry contacts are key, since
it's often hard or impossible to do compelling research in a systems area
without collaboration with industry partners. Do the professor's students work
for companies that excite you? Do the advisor's current students often do
consulting or summer internships in industry?

\subsection{Switching Advisors}

It is not recommended to switch advisors after your second year, but it has
been done. Switching within or immediately after your first year is common and should not
be feared. If you want to switch, just have a candid conversation with your
advisor; be honest and try to leave on peaceful terms.

\subsection{Talking with your Advisor}

The role of your advisor is to provide `advice'. If you disagree
with a decision he/she has made, you should talk to him/her about it. Disagreements
may come in many forms (research direction, summer plans, or even funding).

You should talk to your advisor often, at least once a week. There's a number
of topics you may need to talk about during your \phd:

\begin{itemize}
\item {\it Be explicit about expectations.} How often does your advisor
expect you to publish? Specifically, what about the upcoming {\sc foo-bar}
conference? It's best to get your advisor to be upfront about what he/she
expects than to disappoint.

\item {\it Make meetings count.} Come to meetings
prepared with a plan of what to talk about, questions, etc., and leave with a
clear idea of what you'll do before the next meeting.

\item {\it Paper author order.} A number of groups have strife because
there is no clear policy, managed by the advisor, to control how authors are
ordered on papers. Don't wait until the paper is completed to discuss author
order---discuss it with your advisor ahead of time. Likewise, working on projects
with students on the same graduation schedule as you may be problematic, as you
will both want first authorship on the same paper. Working with junior students 
(when you are a senior student) will result in fewer problems.

\item {\it Theses from large projects.} Large, multi-year projects can
involve an unknown number of graduate students, but usually produce a finite
number of theses. If you get involved in one of these projects, be sure to
discuss what part of the project can be used for your thesis. If that can't be
determined within the first few quarters, you may want to find a new project.
\end{itemize}

\section{\secso{Funding}}

Funding is vital; you should have candid talks with your advisor about it,
so you can ensure support through your entire graduate career with minimum
distraction from your research. Your happiness is often correlated with
funding (and also with progress). Grad students with funding tend to be pretty
happy, but without funding, grad life can be pretty miserable. Most students 
find their thesis topic by working in an area or project that provides
them funding.

Every quarter you need 50\% of support. Funding comes in percent units,
usually 50\% or 25\%. Ideally you are a 50\% Research Assistant ({\sc ra}),
funded under a grant from a government agency ({\sc doe}, {\sc dod}, {\sc nsf})
or a private organization, and working on an interesting topic.

You may be lucky and have a fellowship from Stanford, {\sc nsf}, {\sc ndseg},
Hertz, or some other group that covers you for three to five years. If you're
really lucky, you have two fellowships, one deferrable, like {\sc nsf} and {\sc
ndseg}: defer {\sc nsf} for two years, and get five years of coverage. If it's your
first year and you don't have a fellowship, you should apply for {\sc nsf},
which still open to graduate students in their first and second year. There are
also one year fellowships from companies like Microsoft, Intel, {\sc ati},
{\sc nvidia}, and Yahoo.

If you're not on a grant or fellowship, then you're a Course Assistant ({\sc
ca}). Unless you're going for the Distinction in
Teaching, this can be a bad situation. Being a {\sc ca} requires $\approx 20$
hours a week of work before you've done a single minute of research. Do not be
a long-term {\sc ca}. You will be unable to do enough research 
to keep your advisor happy, and so you continue on as a {\sc ca},
in a negatively reinforcing cycle.  If you have to {\sc ca}, ask the senior
grad students what are the good courses to {\sc ca} that don't require 20 hours
of work. A rule of thumb might be to avoid intensive programming project
courses like Operating Systems.

\subsection{Grants}

Grants are given to your advisor and your advisor employs you as a {\sc ra}. It is
rare, but not unheard of, that a professor has cut a student's funding,
although I believe they are required to give a one quarter notice so that you
have time to scramble to find other funding. This is something that should not
happen to you if you are having candid talks with your advisor about your
progress.

You should know what grant you are on, if it is renewable, and what the
criteria is for renewal. You should also know if your advisor is applying for
other grants from the {\sc nsf}, {\sc dod}, {\sc doe}, or whomever. Try and
help out on the grant writing. If you want to be faculty in the future you
should know how to write a grant, in any case. If your advisor is not writing a
grant, you should be concerned. In this case, ask your advisor if they will
help you write a grant.  You can go to the {\sc nsf} website and look at the calls.

\subsection{50\% vs. 25\% CA}

Avoid a 25\% {\sc ca} (Course Assistantship), as you usually do
nearly the same amount of work of a 50\% {\sc ca} (which fully funds you for
the quarter). As a 25\% {\sc ca}, you'll need to find another source of 25\%
funding (possibly a 25\% {\sc ra} with a professor who is not your advisor).

\subsection{50\% vs. 25\% RA}

Most {\sc ra}ships are 50\% (except in the summer, when they are 90\%, which is
1.8$\times$ your normal salary). Usually you are an {\sc ra} under your
professor. Sometimes people take an {\sc ra} position with someone besides their
professor to have support.  If this happens to you, you can minimize the impact
on your research by finding an {\sc ra} position that is somewhat related to
your research (you can kickstart a collaboration between your advisor and another professor).
Sometimes, an {\sc ra} under another professor may have you writing code for some group 
in the university that needs it.
Students in the {\sc ee} department tend to know a little more about doing
{\sc ra}ships not for your advisor, so try asking one of them.

\subsection{Fellowships}

Fellowships can be a double-edged sword. It's great to have long-term funding
with few or no strings attached, but the lack of quarter-to-quarter
accountability for funding can often lead to drift, especially in the case of a
hands-off advisor. Celebrate, but beware! A fellowship can enable you to work
with the advisor you want to work with, because advisors are usually very
welcoming of a student that comes with funding.

\subsection{Consulting}

You can officially consult for other companies for one day a week. It's a nice
way to make money and may not require too much time on your part.

\subsection{Summer Internships}

Internships can be a great source of extra spending money, a valuable way to
hook into industry research, and a way to get ideas for a thesis. Spend at least 
one summer at an industry research lab. We know of students who took their summer 
internship at an industry research lab (e.g., Microsoft Research) and
turned it into a thesis. Your advisor may have contacts, and can recommend
internships for you. Some good places are: Microsoft Research, Microsoft
Research Asia, Yahoo Research, Intel Research, {\sc hp} Labs, {\sc parc}, or
your advisor's company.

\subsection{Unrestricted Funds}

Grants are reserved for particular purposes and projects. So, you (usually)
can't get money from the {\sc foo} pot unless you work on the {\sc foo}
project. However, your advisor may have some unrestricted funds, and if you
are lucky, they may be able to cover you. Unrestricted funds are often
coveted, so don't rely on this.

\subsection{TGR}

If you're in your last year, you can make your fellowship money stretch father 
by going {\sc tgr}, which reduces the amount of tuition.
However, you can
technically only stay on {\sc tgr} for a year, though this rule is probably
flexible.

\subsection{Graduation Quarter}

If you just need one more quarter, the department allows you to register for a
graduation quarter which costs very little (\$100). If you've saved up some
money and really are on your last quarter, you can take this option. Be
absolutely positive it is your last quarter, because if it isn't, you'll have
to pay a penalty to the department.

\section{\secso{Comps}}

\subsection{Study for Comps}

No matter how much they say you should not study for Comps, you should
really form study groups (to get to know other first years) and study as hard
as you can for Comps so you can get them out of the way. You don't really
want to re-take those undergraduate classes again. Not getting them out of the
way early simply means that next fall, when you're working on research,
you're going to have to take time to study for them again.

Another point of view: Comps are not nearly as important as finding an advisor.
Focus on Comps only if you know who you are working with and are confident you
will have funding and an advisor next quarter. Wasting all your time on Comps
at the expense of research is a terrible idea unless you have an outside
fellowship or you have already convinced a professor to take you.

\subsection{Take the Class}

Also, if you know you have a slight chance on passing a Comp, take the
class if possible in the fall. Often people pass the Comp and can drop the
class even though only about half of the class will have been taught when Comps
begin. If you don't pass the Comp, you can step things up in the class and get
the required A-minus to get out of the comp.

\subsection{Watch the Lectures}

Alternatively, if you don't want to take the class, it is useful to watch the
online lectures. This is especially useful if the person who is writing the
Comp is actually teaching the course that quarter. For example, if you happen
to find out that Prof.~ Koller is writing the {\sc ai} comp, and she is teaching 121
that quarter, it will probably be useful to watch the online version of her
lectures.

http://scpd.stanford.edu/scpd/default.htm

\subsection{Read the Book and Take Practice Exams}

It is entirely possible to pass some Comps by just reading the associated book.
One great example is reading Michael Sipser's short book and passing 
Automata.

\subsection{Form a Study Group}

This is a great way to meet other first years. The 2006--2007 first-years
formed study groups and met weekly to keep each other motivated, and also
posted relevant materials on a wiki:

http://phdcomp.pbwiki.com/

\subsection{Take the Comp Anyways}

Sometimes it helps with the curve if more people sit for a comp. So even if you
don't know very much about the Comp, sitting for it may help your friends, and
vice versa.

\section{\secso{How to do Research}}

An undergraduate degree is about doing well in classes: someone tells you
what to do and how to do it. Graduate school is very different. You are solely
responsible for setting your daily schedule, finding the right project to work
on, and figuring out how to accomplish it. This uncertainty may be difficult to
adjust to and can contribute to a lot of stress. There may be long stretches of
time when it is unclear if your work will be fruitful.

Here are some tips on how to avoid this uncertainty...

\subsection{Read Papers}

It is critical to read {\it lots} of papers in your area, and more importantly, in
your particular research topic. What if you don't like to read papers? Then at
the very least, {\it skim} lots of papers.

\subsection{Talk With People}

Conversations make you smarter. They help you look at problems from different
perspectives, hone your arguments, refine your terminology, explore alternative
solutions, and uncover existing approaches. Chat with professors, other
graduate students (even outside your area), and anyone else who will give you
the time of day. Not only will this improve your research, but more people will
know you and what you are doing. Make sure you are listening at least as much
as you are talking.

Bounce your ideas off other colleagues. If they like your work, great. If not,
ask them why and improve your work based on their feedback.

\subsection{Be Goal Oriented}

When deciding on research projects, be goal oriented. Do not waste your time
working on projects that lead nowhere. Ask yourself, ``Will this project lead to
a paper? Which conference will you submit it to? Is it a chapter in the
thesis?'' Working on a project simply because it's cool is not a
sufficient reason.

\subsection{Set Deliverables and Discuss with your Advisor}

Usually a project will span weeks if not months or years. Often, students
will be distracted and/or discouraged. Set weekly goals. Turn these goals into
deliverables that you can show or discuss with your advisor and colleagues. If
you're going the wrong direction, the early feedback will help you get back on
track.

\subsection{Time Management and Procrastination}
Some (read: many) of us find it challenging to manage our schedules. If you really
need to get work done, consider techniques to avoid time wasting. If you close
your office door (or wear headphones), fewer folks will bother you. If your office
is too loud, consider switching offices (or even working in the campus libraries).


\section{\secso{Quals}}

The milestone exam that tests the ``depth'' of your knowledge in your area is Quals. Some
folks take it their second year; some take it their third. Ideally, you should
pass quals your first time through. Officially, it's 100\% {\sc ok} to not pass
your first time, and then pass your second. It's possible to take quals three
or more times, but your advisor better be loving your research. There have been
students who have left with a Masters after failing twice; other students have
passed on the third try, generally based on if their advisor likes what they
are doing or not.

In general, orals can be ``better'' than written exams: if you know something
about the area, the questioner can tease it out, whereas on a written exam
there is no back and forth. Definitely try your best the first time around
(form study groups).


\section{\secso{The Thesis}}

Theses vary between disciplines, and even advisors. Yours should be modeled after
other theses in your area: grab some examples and discuss the structure with
your advisor. A rule of thumb that has been described (in Graphics or {\sc
hci}) is to have about 3 papers on a single topic, which then become the 3 main
chapters of your thesis. An effective strategy for initiating the process of
completing your thesis has been to present your advisor with an outline of your
proposed thesis. This allows the both of you to discuss what, if anything, may
be missing from your research and can allow you to target your work so that you
can graduate in a timely manner.

\subsection{Committee}

You need to have a committee (see the \phd\ Bible)
or other source for composition and
forms). Make sure your reading committee will have time to read and sign your
thesis. It is better to have committee members who will be able to give timely
feedback about your thesis. You should meet somewhat regularly with your other
committee members to update them about your progress. You don't want them to be
surprised at your defense about the direction you are taking for your thesis!
Better for you to hear that they don't like what you're doing 12 months before
your defense than the day of...

\subsection{Defense}

Ideally your advisor shouldn't let you schedule a defense if you aren't ready,
but there have been students who have not passed the orals, even though their
advisor let them defend. Most people usually defend after having written some
(but not all) of their thesis. Scheduling your orals should probably be done
two to three months in advance. Don't forget to bring food. Catering means one
less thing to worry about.


\subsection{Writing}

Thesis lengths vary from under a hundred pages to as long as several hundred pages. Some
advisors are comfortable with a ``stapled thesis,'' which is merely the papers
you've written with transition paragraphs between them. Other advisors
want the thesis to stand on its own as a treatise on a subject; this
requires some amount of rewriting.


\chapter{Health and Happiness}

\section{\secso{You Are An Asset}}

A graduate student costs more than \$100{\sc k/yr} to fund, including tuition,
stipend, office space, computing equipment, and travel expenses. You are an
enormous investment. A graduate student who is unable to do research (due to
{\sc rsi}, depression, or other physical or mental health issues) is a waste of that
money. Your advisor is likely to accommodate your needs where your well-being is
at stake. Recognize your value, and put your health above all else.

\section{\secso{Mental Health}}

A \phd\ can be a high-stress, low-social interaction environment that can lead
to a variety of mental issues, including anxiety and depression.

A way to be more involved in Stanford life is to take/audit fun
classes such as languages, photography, singing, piano, or art. Having
something to go to that is scheduled every week can provide structure and be a
nice way to make friends. Playing intramural sports (e.g., Basketball, Soccer,
Ultimate Frisbee) or joining one of the larger clubs might also be a good idea:
Taekwondo and Wushu provide a great team experience and has grads and
undergrads. The Stanford Outdoors clubs are a good choice as well. Some grad
students do various kinds of dance: Viennese Ball or the {\sc sbdt}.

Keep in mind that if you're stressed or feeling worried about the uncertainty
of grad school, other students are as well. You're not alone in this. Talk to
other students about it. It will make you feel better. If you feel you need
more help, Counseling And Psychological Services ({\sc caps}) is an arm of Vaden
Health Center designed to help with mental health.

The \phd\ Comic Strip has a forum populated with graduate students
from around the world, which can be a great source of support:

http://www.phdcomics.com/proceedings/

\section{\secso{Physical Health}}

Although we aren't medical doctors, we think that regular exercise is great for
reducing stress (which is pretty high in a \phd\ program) apart from its other
benefits. Take advantage of the Stanford athletic
facilities: apart from the gym, there are classes in tennis, swimming,
climbing, yoga, golf, volleyball. You also happen to be in the Bay Area where
there is great skiing a few hours away, surfing in Santa Cruz, hiking and
climbing in Yosemite, and good weather year round.

You should also eat right.
The free food you find around Gates is probably not the most healthy...

\subsection{RSI}

You probably do a lot of typing, and you may or may not develop {\sc rsi}. Educate
yourself about ergonomics. Consider getting an ergonomic keyboard (Kinesis Maxim or Goldtouch). Some students use Workrave which
monitors typing and mousing behavior and reminds you to take breaks every once
in a while. Students report that this has significantly improved their {\sc rsi}. The
defaults are if you're already experiencing pain (and you should obey the
limits set by the program). If you just want to be preventative, set the
defaults to be a little less strict.

\section{\secso{Making Friends}}

Generally, you will start hanging out with folks in your incoming class, and in
your research lab. If you would like to meet folks outside of {\sc cs}, there
are many opportunities. You can take classes in Music, Foreign Languages,
Physical Education, and Social Dance. One way to meet many folks is to become a
{\sc ca} (Community Assistant) in your graduate housing complex. There are also
outdoors clubs (Snowboarding, Hiking, Running).

Meet more people early on in your career, as you will be busier as time goes on.

\chapter{Department}
\section{\secso{Paperwork and Pitfalls}}
There are several types of paperwork that matter, and they generally deal with
getting paid and getting your \phd\ milestones complete.

When you pass enough Comps, you should file your candidacy form, so you get a
pay increase. If you have some strange advising situation (like your advisor is
not natively {\sc cs} or has a joint appointment with some other department),
make sure it goes through.

Before each quarter, make sure you have funding. If you are {\sc ra}ing, you
should sign a form with your advisor's admin.

When you near the end of your \phd, you will need to fill out a form for
your Reading Committee, and one to schedule your Oral Examination.

Within the last year, you can go {\sc tgr} to save your lab money.

Here are where you can find many of the forms:

http://registrar.stanford.edu/shared/forms.htm\#GradStud

\vspace*{-5pt}
\section{\secso{Lounge}}
\vspace*{-5pt}
There are many elected positions, but some worth mentioning are the lounge
committee members. {\it You}, the student body, are in charge of the lounge.
Make sure there are ping-pong balls and foosballs (money is available from the
department) and that the lounge stays tidy.

\chapter{Life After Grad School}

\section{\secso{Doing Something Else}}

About 10\% of \phd\ students, in our experience, drop out to pursue other
interests, although we don't know the official statistics. They might leave
with a Master's degree or without. There are legitimate reasons to stay or to
go. If you want to be faculty, or if you want to do research in the future, you
need a \phd. However, you may be unable to work with your advisor, unable to
find an advisor to work with, or doing research may not be what you expected or
want. Be aware that there is an opportunity cost to staying in graduate school
since we can make very good salaries working in industry (if you estimate a
salary as \$100{\sc k} and a graduate student stipend as \$30{\sc k}, \$70{\sc
k} is what you're ``giving up'' to be here). There is also a cost to leaving
graduate school. Be clear with yourself about why you are still in the program
or why you want to leave. If you do leave, try not to burn any bridges.

\section{\secso{Becoming Faculty}}

About 40--50\% of Stanford \phd\ graduates go on to faculty careers.  This is a
lower percentage than at other schools, because of Stanford's strong industry
ties.  As with much of grad school, successfully preparing for and landing a
faculty position won't just happen: you need to take initiative.

\subsection{Kinds of Faculty Positions}
Only about 5\% of Stanford \phd s (one or two students in your entering class) will end up
in a job that resembles their
advisor's tenure track position at a top-5 {\sc cs} department. Others will end
up lower on the ranking totem-pole or in other kinds of academic jobs. If you
are interested in academia, you should think about the basis of your interest.
Is it doing research?  Freedom to follow your interests?  Advising students?
Teaching?  Working with smart people? Or just the university environment?

There is a pecking-order of prestige in academia:

\begin{enumerate}
\setlength{\itemsep}{1pt}
\setlength{\parskip}{0pt}
\setlength{\parsep}{0pt}
\item tenure-track faculty at a top research school
\item tenure-track faculty at other research universities
\item teaching-centered academic jobs
\begin{itemize}
\item teaching faculty at a research school (non-tenure)
\item tenure-track faculty at a liberal arts college
\end{itemize}
\item non-faculty research staff
\item adjunct faculty
\end{enumerate}

Higher up in this pecking order, the jobs are more demanding, harder to get,
and pay better.  You also get better students and have an easier time getting
funding the higher you are.  Different positions also emphasize different parts
of the job.  Consider how hard you want to work and what kind of work you want
to do. If you are in doubt about what you want, aim high.  It is generally easy
to switch from a research-centered career to a teaching-centered career, and to
slide down the totem pole within either domain, but it is extremely hard to
move the other way once your career has begun.  Most academia-bound Stanford
\phd s end up at number 2, with a few stars hitting number 1 and a few deciding
to focus on teaching.  Avoid an adjunct career.

\subsection{Preparing}

Whatever sort of academic career you're heading for, you should consider your
time in grad school to be an apprenticeship.  Find opportunities to do the
things faculty do.

\begin{itemize}
\item Help your advisor write grants.  Learn how funding works.
\item Recruit junior grad students or undergrads for your project and act as a
mentor for them. Participate in {\sc curis}, which pairs
undergrads with graduate students during the summer. If you get a good
undergrad and you have a well-defined project, they can be extremely helpful.

\item Submit papers, attend conferences, and give lots of talks.

\item Network with other researchers at conferences and when they come to give
talks at Stanford.

\item Choose {\sc ta}ships where you can help write tests or assignments and
give lectures.

\item Be a teaching fellow and teach a whole class.

\item Earn the department's Distinction in Teaching.
\end{itemize}

There are several specific resources available at Stanford:

\begin{itemize}
\item The Future Faculty Seminar ({\sc cs} 400), offered each spring, is a set of
talks on all aspects of an engineering academic career, mostly given by
Stanford faculty.

\item The Career Development Center offers a series of talks on academic
careers.

\item The Center for Teaching and Learning has lots of talks and resources for
becoming a better teacher.
\end{itemize}

\subsection{Readings}
{\it Tomorrow's Professor}, by Rick Reis. A very good book on what to do while
you're in grad school in order to land a faculty job.


\chapter{References}
\vskip8pt
The Redbook Wiki\\*
\verb+http://theredbook.pbwiki.com+\\*
\\*
Ron Azuma. ``So Long, and thanks for the Phd!''\\*
\verb+http://www.cs.unc.edu/~azuma/hitch4.html+\\*
\\*
Jason Hong. ``Grad School Advice''\\*
\verb+http://www.cs.cmu.edu/~jasonh/advice.html+\\*
\\*
\phd\ Bible\\*
\verb+http://www-cs-students.stanford.edu/phd/+ 


\chapter{About this Document}
\vskip8pt
In Spring of 2007, a panel of visiting faculty
was conducting an every 5--6 year evaluation of the Stanford {\sc cs} department. It
was mentioned that Purdue students distribute an advice
manual, called the Blue Book, which is passed from generation to
generation of students. Not to be outdone, we created the Redbook.

\vspace*{-5pt}
\section{\secso{Philosophy and Style}}
\vspace*{-5pt}

The Redbook should be concise to encourage complete reading and facilitate future updates. It should offer advice instead of listing facts about the program, because those facts can be gotten elsewhere.

\vspace*{-5pt}
\section{\secso{Maintenance}}
\vspace*{-5pt}
Please maintain the Redbook: add to it, edit it, correct it, and distribute it
to first year students in the Fall. The 2007 printing was supported by the
Computer Forum although the content was created independently.

Please append a dedication and add to the acknowledgments page.

{\footnotesize
The Redbook has an email. Use it to communicate with publishers and
advertisers.\\*
\verb+stanfordredbook@gmail.com / ask a student+

This edition was printed by CafePress.com:\\*
\verb+http://www.cafepress.com/theredbook+.\\*
\verb+stanfordredbook@gmail.com / ask a student+

Find source materials at\\*
\verb+http://www.austenmcdonald.com/projects/redbook+}

\vfill\eject
\thispagestyle{empty}
\mbox{}
\vfill\eject
\thispagestyle{empty}
\mbox{}
\vfill\eject
\thispagestyle{empty}
\mbox{}
\vfill\eject
\thispagestyle{empty}
\mbox{}
\vfill\eject
\thispagestyle{empty}
\mbox{}
\vfill\eject
\thispagestyle{empty}
\mbox{}
\vfill\eject
\thispagestyle{empty}
\mbox{}
\vfill\eject
\thispagestyle{empty}
\mbox{}
\vfill\eject

\end{document}
